% !TEX root = ../../hw05.tex
\begin{proof}[Solution]
  A fixed point is a point where $F(x)$ and $x$ have the same value. To
  find one, consider their difference,
  \[
    G(x)=F(x)-x.
  \]
  This function is continuous on $[a,b]$. Since $F$ maps $[a,b]$ into
  itself, $F(a)\ge a$ and $F(b)\le b$. Hence
  \[
    G(a)\ge0,\qquad G(b)\le0.
  \]
  If either value is zero, the corresponding endpoint is a fixed point.
  Otherwise $G(a)>0>G(b)$, so the intermediate value theorem gives a
  $c\in(a,b)$ with $G(c)=0$. Thus $F(c)=c$. This also covers a
  one-point interval, since then $F(a)=a$.

  Now fix $n\in\N$ and define
  \[
    F_n(x)=\frac{1}{1+x^n},\qquad x\in[0,1].
  \]
  The denominator is positive on this interval, so $F_n$ is continuous.
  Also $0\le x^n\le1$, which gives
  \[
    \frac12\le F_n(x)\le1.
  \]
  Therefore $F_n$ maps $[0,1]$ into itself. By the result just proved,
  it has a fixed point $c\in[0,1]$. Neither endpoint is fixed, since
  $F_n(0)=1$ and $F_n(1)=1/2$. Consequently $c\in(0,1)$ and
  \[
    c=\frac{1}{1+c^n}.
  \]
  \autoref{fig:hw05:interval-fixed-point} shows the intersection when
  $n=1$.

  \begin{marginfigure}
    \centering
    % !TEX root = ../../problems/hw05.tex
% Shared physical dimensions for interval and number-line diagrams.
\tikzset{
  interval diagram/.style={
    x=1cm,y=6mm,font=\footnotesize,
    color=black,line width=0.5pt,line cap=round,line join=round,
    >={Stealth[length=4pt,width=3pt]}
  },
  interval endpoint/.style={
    circle,draw=black,line width=0.5pt,
    minimum size=4pt,inner sep=0pt,outer sep=0pt
  },
  interval open/.style={interval endpoint,fill=white},
  interval closed/.style={interval endpoint,fill=black},
  interval label/.style={inner sep=1pt}
}

\begin{tikzpicture}[interval diagram,x=28mm,y=28mm]
  \draw[->] (-0.05,0) -- (1.13,0) node[right=4pt] {$x$};
  \draw[->] (0,-0.05) -- (0,1.14) node[above=4pt] {$y$};
  \node[below left=5pt] at (0,0) {$0$};
  \draw[gray] (1,-0.025) -- (1,0.025);
  \draw[gray] (-0.025,1) -- (0.025,1);
  \node[below=6pt] at (1,0) {$1$};
  \node[left=5pt] at (0,1) {$1$};

  \draw[gray,densely dashed] (0,0) -- (1.04,1.04);
  \node at (0.68,1.09) {$y=x$};
  \draw[domain=0:1,samples=80] plot (\x,{1/(1+\x)});
  \node at (0.92,0.76) {$F_1$};
  \fill (0,1) circle[radius=1.6pt];
  \fill (1,0.5) circle[radius=1.6pt];

  \draw[gray,densely dotted] (0.618034,0) -- (0.618034,0.618034)
    -- (0,0.618034);
  \fill (0.618034,0.618034) circle[radius=1.6pt];
  \node[below=6pt] at (0.618034,0) {$c$};
  \node[left=5pt] at (0,0.618034) {$c$};
\end{tikzpicture}

    \caption[A fixed point as a graph intersection.]{%
      For $n=1$, the graph $y=1/(1+x)$ crosses $y=x$
      inside the unit square. At the crossing, the input equals its output.}
    \label{fig:hw05:interval-fixed-point}
  \end{marginfigure}

  When $n=1$, multiplying by $1+c>0$ gives
  \begin{align*}
    c=\frac{1}{1+c}
      &\iff c(1+c)=1\\
      &\iff c^2+c-1=0\\
      &\iff c=\frac{-1+\sqrt5}{2}
        \quad\text{or}\quad c=\frac{-1-\sqrt5}{2}.
  \end{align*}
  Only the first value belongs to $(0,1)$. Hence the required value is
  $c=(\sqrt5-1)/2$.
\end{proof}
